% block: krylov-classification-heading page:85 type:chapter
\clearpage\section{\BT{Classification and Theory of Krylov Subspace Methods}{Krylov 子空间方法的分类与理论}}\label{chap:krylov-classification}
\setcounter{footnote}{0}
% block: krylov-classification-abstract page:85 type:ordinary-paragraph
\noindent\textbf{\BT{Abstract}{摘要}}\quad
\BT{Krylov subspace methods are roughly classified into three groups: ones for Hermitian linear systems, for complex symmetric linear systems, and for non-Hermitian linear systems. Non-Hermitian linear systems include complex symmetric linear systems since a complex symmetric matrix is non-Hermitian and symmetric. Krylov subspace methods for complex symmetric linear systems use the symmetry of the coefficient matrix, leading to more efficient Krylov subspace methods than ones for non-Hermitian linear systems. This chapter also presents preconditioning techniques to boost the speed of convergence of Krylov subspace methods.}{Krylov 子空间方法大致分为三类: 用于 Hermitian 线性方程组的方法, 用于复对称线性方程组的方法, 以及用于非 Hermitian 线性方程组的方法. 非 Hermitian 线性方程组包含复对称线性方程组, 因为复对称矩阵是非 Hermitian 的且是对称的. 复对称线性方程组的 Krylov 子空间方法利用系数矩阵的对称性, 因而比非 Hermitian 线性方程组的 Krylov 子空间方法更高效. 本章还介绍用于提高 Krylov 子空间方法收敛速度的预条件技术.}
